\documentclass[12pt]{article}

\usepackage[a4paper, top=2cm, bottom=2cm, left=2.5cm, right=2.5cm]{geometry}

\usepackage{amsmath, amssymb}
\usepackage{tikz}
\usetikzlibrary{arrows.meta,positioning,fit,calc}
\usepackage{multirow}
\usepackage{array}
\usepackage{siunitx}
\usepackage{enumitem}
\usepackage{float}
\usepackage{fontspec}
\usetikzlibrary{decorations.pathmorphing}
\usepackage{pdflscape}
\usepackage{tabularx}
\usepackage{booktabs}
\usepackage{graphicx}
\usepackage{listings}
\usepackage[table]{xcolor}
\usepackage{hyperref}
\usepackage{bookmark}
\usepackage{pdfpages}

% ---------------- FONTS ----------------
\setmainfont[
    UprightFont = Handlee-Regular.ttf,
    AutoFakeBold = 2.5
]{Handlee}

\newfontfamily\headingfont{PT.ttf}
\setmonofont{DejaVu Sans Mono}

% ---------------- HANDWRITTEN MATH ----------------
\usepackage{mathastext}
\Mathastext
\everymath{\displaystyle}

% ---------------- GENERAL ----------------
\setlength{\parindent}{0pt}
\pagenumbering{gobble}
\linespread{1.2}

\newcommand{\mychapter}[2]{%
    \phantomsection
    \hypertarget{#2}{}%
    \bookmark[level=0,dest=#2]{#1}%
    \begin{center}
        {\headingfont\Huge #1}
        \par
        \vspace{4pt}
        \hrule
    \end{center}
}

\newcommand{\mysection}[3]{%
    \phantomsection
    \hypertarget{#3}{}%
    \bookmark[level=1,dest=#3]{#1\space #2}%
    \newpage
    \noindent{\headingfont\Large #1\space #2\par}
    \vspace{4pt}
    \hrule
    \vspace{15pt}
}

\newcommand{\mysubsection}[2]{%
    \vspace{5px}
    \noindent{\headingfont\large #1\space #2\par}
    \vspace{4pt}
    \hrule
    \vspace{15pt}
}

\begin{document}

\mychapter{Chapter 3 : Properties of Pure Substance}{5f4y5fdh}
\begin{center}
\begin{tabularx}{\textwidth}{|>{\centering\arraybackslash}p{2.0cm}|X|}
\hline
\textbf{Section} & \textbf{Core Revision Focus} \\
\hline

3.1 &
Pure substance definition, fixed chemical composition, purity criterion, single/multiple components and phases
\\
\hline

3.2 &
Solid/liquid/gas phases, molecular arrangement and motion, intermolecular forces, phase characteristics
\\
\hline

3.3 &
Phase-change processes, compressed/saturated/superheated states, saturation temperature/pressure, latent heat, phase-change sequence
\\
\hline

3.4 &
$T$-$v$, $P$-$v$, $P$-$T$ diagrams, saturation dome, critical point, triple point, sublimation, phase regions, $P$-$v$-$T$ surface
\\
\hline

3.5 &
Property tables, saturated-liquid/vapor properties, quality, mixture-property relations, superheated vapor, compressed liquid, state identification
\\
\hline

3.6 &
Ideal-gas equation $Pv=RT$, specific gas constant, molar form, two-state relation, constant-$P/V/T$ cases, density/specific volume, ideal-gas applicability
\\
\hline

3.7 &
Compressibility factor $Z$, deviation from ideal-gas behavior, reduced properties, principle of corresponding states, generalized compressibility chart
\\
\hline

3.8 &
van der Waals, Beattie--Bridgeman, Benedict--Webb--Rubin, Strobridge and virial equations, real-gas modeling, applicability and accuracy
\\
\hline

\end{tabularx}
\end{center}
% ============================================================
\mysection{3.1}{Pure Substance}{3.1}

\textbf{\\ Definition}

A \textbf{pure substance} has a \textbf{fixed chemical composition throughout}.

\begin{equation*}
\boxed{\text{Uniform chemical composition throughout}}
\end{equation*}

\textbf{\\ Key Criterion}

A substance may contain one component, several components, or multiple phases and still be pure provided the composition is uniform.

\begin{equation*}
\boxed{
\text{Same composition in all regions/phases}
\Rightarrow
\text{Pure substance}
}
\end{equation*}

\begin{equation*}
\boxed{
\text{Different composition in different regions/phases}
\Rightarrow
\text{Not a pure substance}
}
\end{equation*}

\begin{center}
\begin{tabular}{|c|c|}
\hline
\textbf{Mixture} & \textbf{Pure Substance?} \\
\hline
Ice + liquid water & Yes \\
\hline
Liquid air + gaseous air & No \\
\hline
Oil + water & No \\
\hline
\end{tabular}
\end{center}

\textbf{\\ Exam Point}

\begin{itemize}[leftmargin=*]
    \item Number of components or phases alone does \textbf{not} determine purity.
    \item Multiple phases are permissible if all phases have the same chemical composition.
\end{itemize}

% ============================================================
\mysection{3.2}{Phases of a Pure Substance}{3.2}

\textbf{\\ Principal Phases}

\[
\boxed{\text{Solid,\quad Liquid,\quad Gas}}
\]

\textbf{\\ Definition of a Phase}

A phase has:
\begin{itemize}[leftmargin=*]
    \item a distinct molecular arrangement,
    \item a homogeneous structure, and
    \item identifiable boundaries separating it from other phases.
\end{itemize}

\textbf{\\ Molecular Comparison}

\begin{center}
\begin{tabular}{|c|c|c|c|}
\hline
\textbf{Property} & \textbf{Solid} & \textbf{Liquid} & \textbf{Gas} \\
\hline
Molecular arrangement & Ordered lattice & Less ordered & No order \\
\hline
Position & Relatively fixed & Not fixed & Widely separated \\
\hline
Motion & Oscillation & Rotation + translation & Random \\
\hline
Intermolecular force & Strongest & Intermediate & Weakest \\
\hline
Molecular energy & Lowest & Intermediate & Highest \\
\hline
\end{tabular}
\end{center}

\textbf{\\ Phase Transformation}

\begin{equation*}
\boxed{
T\uparrow
\Rightarrow
\text{Molecular velocity}\uparrow
\Rightarrow
\text{Intermolecular forces can be overcome}
}
\end{equation*}

\textbf{\\ Multiple Phases}

The same substance may possess different phases within one principal phase because of different molecular structures.

Examples from the notes: graphite/diamond, multiple liquid phases of helium, multiple solid phases of iron, and multiple high-pressure ice phases.

% ============================================================
\mysection{3.3}{Phase-Change Processes of Pure Substances}{3.3}

\textbf{\\ Phase-Change Sequence}

For constant-pressure heating of a pure substance:

\begin{equation*}
\boxed{
\text{Compressed liquid}
\rightarrow
\text{Saturated liquid}
\rightarrow
\text{Saturated mixture}
\rightarrow
\text{Saturated vapor}
\rightarrow
\text{Superheated vapor}
}
\end{equation*}

\textbf{\\ Compressed Liquid / Subcooled Liquid}

Liquid \textbf{not about to vaporize}.

\textbf{Saturated Liquid}: liquid \textbf{about to vaporize}.

\[
\boxed{
\text{Compressed liquid}
\rightarrow
\text{Saturated liquid}
\rightarrow
\text{Vaporization}
}
\]

\textbf{\\ Saturated Liquid--Vapor Mixture}

During boiling:
\begin{itemize}[leftmargin=*]
    \item liquid and vapor coexist in equilibrium,
    \item heat transfer continues,
    \item temperature remains constant at constant pressure,
    \item liquid decreases while vapor increases,
    \item specific volume increases substantially.
\end{itemize}

\[
\boxed{
\text{Saturated liquid}
\rightarrow
\text{Saturated mixture}
\rightarrow
\text{Saturated vapor}
}
\]

\textbf{\\ Saturated Vapor}

Vapor \textbf{about to condense}; any heat loss causes condensation.

\textbf{\\ Superheated Vapor}

Vapor \textbf{not about to condense}.

At a specified pressure:
\[
\boxed{T>T_{\mathrm{sat}}\Rightarrow\text{superheated vapor}}
\]

\textbf{\\ Constant-Pressure Phase Change}

During vaporization from saturated liquid to saturated vapor:

\begin{equation*}
\boxed{P=\text{constant},\qquad T=\text{constant}}
\end{equation*}

For water at $1~\mathrm{atm}$:
\[
\boxed{T_{\mathrm{sat}}\approx100^\circ\mathrm{C}}
\]

\textbf{\\ Water and Steam}

Both refer to the same substance:

\[
\boxed{\mathrm{H_2O}}
\]

\textbf{\\ Saturation Temperature and Pressure}

\begin{equation*}
\boxed{
T_{\mathrm{sat}}
=
\text{temperature at which phase change occurs at a specified pressure}
}
\end{equation*}

\begin{equation*}
\boxed{
P_{\mathrm{sat}}
=
\text{pressure at which phase change occurs at a specified temperature}
}
\end{equation*}

\[
\boxed{P_{\mathrm{sat}}=f(T_{\mathrm{sat}})}
\]

\[
\boxed{
P_{\mathrm{sat}}\uparrow
\Rightarrow
T_{\mathrm{sat}}\uparrow
}
\]

For water:
\[
P_{\mathrm{sat}}=101.325~\mathrm{kPa}
\Longleftrightarrow
T_{\mathrm{sat}}=99.97^\circ\mathrm{C}
\]

\textbf{\\ Latent Heat}

Energy absorbed or released during phase change.

\textbf{Latent heat of fusion}
\[
\boxed{h_{\mathrm{fusion}}=333.7~\mathrm{kJ/kg}}
\]
for water at $1~\mathrm{atm}$.

\[
\boxed{\text{Melting: energy absorbed}}
\]
\[
\boxed{\text{Freezing: energy released}}
\]

\textbf{Latent heat of vaporization}
\[
\boxed{h_{\mathrm{vap}}=2256.5~\mathrm{kJ/kg}}
\]
for water at $1~\mathrm{atm}$.

\[
\boxed{\text{Vaporization: energy absorbed}}
\]
\[
\boxed{\text{Condensation: energy released}}
\]

\textbf{\\ Practical Consequences}

\begin{itemize}[leftmargin=*]
    \item Pressure increase $\Rightarrow$ boiling temperature increases.
    \item Elevation increase $\Rightarrow$ atmospheric pressure decreases $\Rightarrow$ boiling temperature decreases.
    \item Higher heat-transfer rate $\Rightarrow$ higher vaporization rate.
    \item Vacuum cooling uses evaporation and latent heat removal to lower product temperature.
    \item Vacuum freezing occurs when
    \[
    P<0.61~\mathrm{kPa}
    \]
    for water at $0^\circ\mathrm{C}$.
\end{itemize}


% ============================================================
\mysection{3.4}{Property Diagrams for Phase-Change Processes}{3.4}

\textbf{\\ Important Diagrams}

\[
\boxed{T-v,\qquad P-v,\qquad P-T}
\]

The complete three-dimensional representation is the \textbf{$P$-$v$-$T$ surface}.

% ------------------------------------------------------------
\textbf{\\ $T$-$v$ Diagram}

As pressure increases:
\begin{itemize}[leftmargin=*]
    \item $T_{\mathrm{sat}}$ increases,
    \item $v_f$ increases,
    \item $v_g$ decreases,
    \item distance between saturated-liquid and saturated-vapor states decreases.
\end{itemize}

\textbf{Critical Point}

At the critical point, saturated liquid and saturated vapor become identical.

Critical properties:
\[
\boxed{T_{\mathrm{cr}},\quad P_{\mathrm{cr}},\quad v_{\mathrm{cr}}}
\]

\[
\boxed{
\text{Saturated liquid line}
+
\text{Saturated vapor line}
\rightarrow
\text{Saturation dome}
}
\]

\begin{center}
\begin{tabular}{|c|c|}
\hline
\textbf{Region} & \textbf{Phase} \\
\hline
Left of saturated-liquid line & Compressed liquid \\
\hline
Under saturation dome & Saturated liquid--vapor mixture \\
\hline
Right of saturated-vapor line & Superheated vapor \\
\hline
\end{tabular}
\end{center}

The left and right regions are \textbf{single-phase regions}; the region under the dome is a \textbf{two-phase equilibrium region} or \textbf{wet region}.

\textbf{Supercritical Region}

\[
\boxed{P>P_{\mathrm{cr}}}
\]

No distinct liquid--vapor phase change occurs; properties vary continuously.

% ------------------------------------------------------------
\textbf{\\ $P$-$v$ Diagram}

The basic regions are the same as in the $T$-$v$ diagram.

During vaporization:
\[
\boxed{
P=\text{constant},\qquad T=\text{constant},\qquad v\uparrow
}
\]

Constant-temperature lines have a \textbf{downward trend}.

% ------------------------------------------------------------
\textbf{\\ Solid Phase and Freezing}

Most substances:
\[
\boxed{\text{contract during freezing}}
\]

Water:
\[
\boxed{\text{expands during freezing}}
\]

On a $P$-$T$ diagram:

\begin{center}
\begin{tabular}{|c|c|}
\hline
\textbf{Behavior} & \textbf{Melting-Line Slope} \\
\hline
Contracts on freezing & Positive \\
\hline
Expands on freezing & Negative \\
\hline
\end{tabular}
\end{center}

% ------------------------------------------------------------
\textbf{\\ Triple-Phase Equilibrium}

\[
\boxed{\text{Solid}+\text{Liquid}+\text{Vapor}}
\]

On $P$-$v$ and $T$-$v$ diagrams: \textbf{triple line}.

On $P$-$T$ diagram: \textbf{triple point}.

For water:
\[
\boxed{T_{\mathrm{tp}}=0.01^\circ\mathrm{C}}
\]
\[
\boxed{P_{\mathrm{tp}}=0.6117~\mathrm{kPa}}
\]

Thus:
\[
\boxed{
T=0.01^\circ\mathrm{C},
\qquad
P=0.6117~\mathrm{kPa}
}
\]

\begin{center}
\begin{tabular}{|c|c|}
\hline
$P-v$ / $T-v$ & Triple line \\
\hline
$P-T$ & Triple point \\
\hline
\end{tabular}
\end{center}

% ------------------------------------------------------------
\textbf{\\ Sublimation}

Direct:
\[
\boxed{\text{Solid}\rightarrow\text{Vapor}}
\]

Without passing through the liquid phase.

Common example: \textbf{dry ice} ($CO_2$).

Sublimation occurs at pressures below the triple-point pressure.

% ------------------------------------------------------------
\textbf{\\ $P$-$T$ Diagram / Phase Diagram}

\begin{center}
\begin{tabular}{|c|c|}
\hline
\textbf{Line} & \textbf{Phases Separated}  \\
\hline
Sublimation line & Solid--Vapor  \\
\hline
Melting line & Solid--Liquid  \\
\hline
Vaporization line & Liquid--Vapor  \\
\hline
\end{tabular}
\end{center}

\[
\boxed{
\text{Sublimation line}
\cap
\text{Melting line}
\cap
\text{Vaporization line}
=
\text{Triple point}
}
\]

\[
\boxed{\text{Vaporization line ends at the critical point}}
\]

Beyond the critical point there is no distinct liquid--vapor boundary.

% ------------------------------------------------------------
\textbf{\\ $P$-$v$-$T$ Surface}

\[
\boxed{z=f(x,y)}
\]

For a simple compressible substance, \textbf{any two independent intensive properties} fix the state.

For the $P$-$v$-$T$ surface:
\begin{itemize}[leftmargin=*]
    \item $T$ and $v$ may be independent,
    \item $P$ is then dependent.
\end{itemize}

Every point on the surface represents an \textbf{equilibrium state}; a quasi-equilibrium process traces a sequence of such states.

\textbf{\\ Diagram Relationships}

\begin{center}
\begin{tabular}{|c|c|}
\hline
\textbf{Diagram} & \textbf{Relation to $P$-$v$-$T$ Surface} \\
\hline
$P-v$ & Projection on $P-v$ plane \\
\hline
$T-v$ & Bird's-eye view \\
\hline
$P-T$ & Projection on $P-T$ plane \\
\hline
\end{tabular}
\end{center}

\textbf{\\ Critical Point vs Triple Point}

\begin{center}
\begin{tabular}{|c|c|c|}
\hline
& \textbf{Critical Point} & \textbf{Triple Point} \\
\hline
Phases & Liquid + vapor become identical & Solid + liquid + vapor coexist \\
\hline
Location & End of vaporization line & Intersection of 3 phase lines \\
\hline
$P-T$ & Point & Point \\
\hline
$T-v$ & Point & Line \\
\hline
$P-v$ & Point & Line \\
\hline
\end{tabular}
\end{center}

% ============================================================
\mysection{3.5}{Property Tables}{3.5}

\textbf{\\ Purpose}

Property tables are used because relationships among thermodynamic properties are often too complex for simple equations.

Typical regions:
\begin{itemize}[leftmargin=*]
    \item saturated liquid / saturated vapor,
    \item saturated liquid--vapor mixture,
    \item superheated vapor,
    \item compressed liquid.
\end{itemize}

% ------------------------------------------------------------
\textbf{\\ Enthalpy}

\[
\boxed{h=u+Pv}
\]

\[
\boxed{H=U+PV}
\]

Therefore:
\[
\boxed{u=h-Pv}
\]

% ------------------------------------------------------------
\textbf{\\ Saturation-Table Notation}

\begin{itemize}[leftmargin=*]
    \item $f$ $\rightarrow$ saturated liquid.
    \item $g$ $\rightarrow$ saturated vapor.
    \item $fg$ $\rightarrow$ vapor minus liquid.
\end{itemize}

\[
\boxed{v_{fg}=v_g-v_f}
\]

\[
\boxed{u_{fg}=u_g-u_f}
\]

\[
\boxed{h_{fg}=h_g-h_f}
\]

$h_{fg}$ is the \textbf{enthalpy of vaporization} and decreases with increasing temperature/pressure, becoming zero at the critical point.

% ------------------------------------------------------------
\textbf{\\ Quality}

Quality $x$ is the vapor mass fraction of a \textbf{saturated liquid--vapor mixture}:

\[
\boxed{x=\frac{m_g}{m_t}}
\]

\[
m_t=m_f+m_g
\]

\[
\boxed{0\leq x\leq1}
\]

\[
x=0\Rightarrow\text{saturated liquid}
\]

\[
x=1\Rightarrow\text{saturated vapor}
\]

\[
\boxed{\text{Quality has meaning only for saturated liquid--vapor mixtures}}
\]

\textbf{\\ Mixture Property Relations}

\[
v=(1-x)v_f+xv_g
\]

\[
\boxed{v=v_f+xv_{fg}}
\]

\[
\boxed{x=\frac{v-v_f}{v_{fg}}}
\]

\[
\boxed{u=u_f+xu_{fg}}
\]

\[
\boxed{h=h_f+xh_{fg}}
\]

General form:
\[
\boxed{y=y_f+xy_{fg}}
\qquad
y=v,u,h
\]

Hence:
\[
\boxed{y_f\leq y\leq y_g}
\]

On $P-v$ or $T-v$ diagrams, quality is the distance ratio:

\[
\boxed{
x=
\frac{\text{distance from saturated liquid to actual state}}
{\text{distance from saturated liquid to saturated vapor}}
}
\]

Low $x$ $\rightarrow$ near saturated liquid; high $x$ $\rightarrow$ near saturated vapor.

% ------------------------------------------------------------
\textbf{\\ Superheated Vapor}

Single-phase vapor region.

\[
\boxed{P<P_{\mathrm{sat}}\quad\text{at a given }T}
\]

\[
\boxed{T>T_{\mathrm{sat}}\quad\text{at a given }P}
\]

\[
\boxed{v>v_g,\qquad u>u_g,\qquad h>h_g}
\]

Quality is \textbf{not defined}.

% ------------------------------------------------------------
\textbf{\\ Compressed Liquid}

When compressed-liquid data are unavailable, use the saturated-liquid approximation at the same temperature:

\[
\boxed{v\approx v_f@T}
\]

\[
\boxed{u\approx u_f@T}
\]

\[
\boxed{h\approx h_f@T}
\]

The approximation is generally good for $v$ and $u$; $h$ is more pressure-sensitive.

Improved enthalpy approximation:

\[
\boxed{
h\approx h_f@T+
v_f@T
\left(P-P_{\mathrm{sat}}@T\right)
}
\]

Characteristics:
\[
\boxed{P>P_{\mathrm{sat}}\quad\text{at a given }T}
\]

\[
\boxed{T<T_{\mathrm{sat}}\quad\text{at a given }P}
\]

\[
\boxed{v<v_f,\qquad u<u_f,\qquad h<h_f}
\]

Quality is \textbf{not defined}.

% ------------------------------------------------------------
\textbf{\\ State Identification}

Given $P,T$:

\[
\boxed{
T<T_{\mathrm{sat}}@P
\Rightarrow
\text{compressed liquid}
}
\]

\[
\boxed{
T=T_{\mathrm{sat}}@P
\Rightarrow
\text{saturated state}
}
\]

\[
\boxed{
T>T_{\mathrm{sat}}@P
\Rightarrow
\text{superheated vapor}
}
\]

Alternatively, given $T,P_{\mathrm{sat}}$:

\[
\boxed{
P>P_{\mathrm{sat}}@T
\Rightarrow
\text{compressed liquid}
}
\]

\[
\boxed{
P=P_{\mathrm{sat}}@T
\Rightarrow
\text{saturated state}
}
\]

\[
\boxed{
P<P_{\mathrm{sat}}@T
\Rightarrow
\text{superheated vapor}
}
\]

Using a property $y=v,u,h$:

\[
\boxed{
y<y_f\Rightarrow\text{compressed liquid}
}
\]

\[
\boxed{
y_f\leq y\leq y_g
\Rightarrow
\text{saturated liquid--vapor mixture}
}
\]

\[
\boxed{
y>y_g\Rightarrow\text{superheated vapor}
}
\]

% ------------------------------------------------------------
\textbf{\\ Reference State}

Absolute values of $u,h,s$ depend on the selected reference state; thermodynamic calculations generally use \textbf{property changes}.

\[
\boxed{
\text{Use one consistent set of tables or charts}
}
\]

Different table sets may have different absolute property values but give consistent property changes.

% ------------------------------------------------------------
\textbf{\\ Table Selection}

\begin{center}
\begin{tabular}{|c|c|}
\hline
\textbf{Known Information / Region} & \textbf{Table} \\
\hline
Saturated state, $T$ given & $T$-based saturation table \\
\hline
Saturated state, $P$ given & $P$-based saturation table \\
\hline
Superheated vapor & Superheated-vapor table \\
\hline
Compressed liquid & Compressed-liquid table \\
\hline
Saturated solid--vapor mixture & Saturated solid--vapor table \\
\hline
\end{tabular}
\end{center}

% ============================================================
\mysection{3.6}{The Ideal-Gas Equation of State}{3.6}

\textbf{\\ Equation of State}

An equation of state relates thermodynamic properties at equilibrium states; the basic form used here relates $P$, $T$, and $v$.

\textbf{\\ Experimental Basis}

Boyle's law:
\[
\boxed{P\propto\frac{1}{V}}
\qquad
(T=\text{constant})
\]

Charles/Gay-Lussac:
\[
\boxed{V\propto T}
\qquad
(P=\text{constant})
\]

% ------------------------------------------------------------
\textbf{\\ Ideal-Gas Equation}

\[
\boxed{Pv=RT}
\]

where:
\begin{itemize}[leftmargin=*]
    \item $P$ = absolute pressure,
    \item $T$ = absolute temperature,
    \item $v$ = specific volume,
    \item $R$ = specific gas constant.
\end{itemize}

\[
\boxed{R=\frac{R_u}{M}}
\]

\[
\boxed{m=MN}
\]

$R_u$ = universal gas constant; $M$ = molar mass; $N$ = number of moles.

% ------------------------------------------------------------
\textbf{\\ Alternative Forms}

\[
\boxed{PV=mRT}
\]

\[
\boxed{PV=NR_uT}
\]

Molar specific volume:
\[
\boxed{\bar v=\frac{V}{N}}
\]

\[
\boxed{P\bar v=R_uT}
\]

Use bars for molar properties; unbarred quantities denote mass-based properties.

% ------------------------------------------------------------
\textbf{\\ Two-State Relation}

For fixed mass of ideal gas:

\[
\boxed{
\frac{P_1V_1}{T_1}
=
\frac{P_2V_2}{T_2}
}
\]

% ------------------------------------------------------------
\textbf{\\ Special Cases}

\textbf{Constant volume}
\[
V=\text{constant}
\]

\[
\boxed{\frac{P}{T}=\text{constant}}
\]

\[
\boxed{\frac{P_1}{T_1}=\frac{P_2}{T_2}}
\]

\textbf{Constant pressure}
\[
P=\text{constant}
\]

\[
\boxed{\frac{V}{T}=\text{constant}}
\]

\[
\boxed{\frac{V_1}{T_1}=\frac{V_2}{T_2}}
\]

\textbf{Constant temperature}
\[
T=\text{constant}
\]

\[
\boxed{PV=\text{constant}}
\]

\[
\boxed{P_1V_1=P_2V_2}
\]

% ------------------------------------------------------------
\textbf{\\ Ideal-Gas Applicability}

Ideal gas = imaginary substance obeying:
\[
\boxed{Pv=RT}
\]

Real gases approach ideal behavior at \textbf{low density}, generally corresponding to:
\begin{itemize}[leftmargin=*]
    \item low pressure,
    \item high temperature.
\end{itemize}

Ideal-gas model is most appropriate at relatively low pressure and high temperature.

Dense gases, steam, and refrigerant vapor may require property tables.

% ------------------------------------------------------------
\textbf{\\ Unit Requirements}

\textbf{Absolute pressure}
\[
\boxed{
P_{\mathrm{abs}}
=
P_{\mathrm{gage}}+P_{\mathrm{atm}}
}
\]

\textbf{Absolute temperature}
\[
\boxed{
T(\mathrm{K})
=
T(^\circ\mathrm{C})+273.15
}
\]

% ------------------------------------------------------------
\textbf{\\ Density and Specific Volume}

Since:
\[
\boxed{v=\frac{1}{\rho}}
\]

\[
\boxed{P=\rho RT}
\]

\[
\boxed{\rho=\frac{P}{RT}}
\]

\[
\boxed{v=\frac{RT}{P}}
\]

Therefore, for fixed $R$:
\[
\boxed{
\rho\uparrow\text{ with }P,
\qquad
\rho\downarrow\text{ with }T
}
\]

\[
\boxed{
v\uparrow\text{ with }T,
\qquad
v\downarrow\text{ with }P
}
\]

% ============================================================
\mysection{3.7}{Compressibility Factor---A Measure of Deviation from Ideal-Gas Behavior}{3.7}

\textbf{\\ Compressibility Factor}

\[
\boxed{Z=\frac{Pv}{RT}}
\]

Thus:
\[
\boxed{Pv=ZRT}
\]

For an ideal gas:
\[
\boxed{Z=1}
\]

Physical interpretation:
\[
\boxed{
Z=\frac{v_{\mathrm{actual}}}{v_{\mathrm{ideal}}}
}
\]

where:
\[
\boxed{
v_{\mathrm{ideal}}=\frac{RT}{P}
}
\]

Therefore:
\[
\boxed{Z=1\Rightarrow\text{ideal-gas behavior}}
\]

\[
\boxed{Z>1\Rightarrow v_{\mathrm{actual}}>v_{\mathrm{ideal}}}
\]

\[
\boxed{Z<1\Rightarrow v_{\mathrm{actual}}<v_{\mathrm{ideal}}}
\]

The farther $Z$ is from unity, the larger the deviation from ideal behavior.

% ------------------------------------------------------------
\textbf{\\ Reduced Properties}

\[
\boxed{
P_R=\frac{P}{P_{\mathrm{cr}}}
}
\]

\[
\boxed{
T_R=\frac{T}{T_{\mathrm{cr}}}
}
\]

Reduced properties are dimensionless.

\textbf{Principle of corresponding states}:

\[
\boxed{
Z=f(P_R,T_R)
}
\]

Different gases have approximately the same $Z$ at the same $P_R$ and $T_R$.

% ------------------------------------------------------------
\textbf{\\ Generalized Compressibility Chart}

The chart gives $Z$ from $P_R$ and $T_R$.

Once $Z$ is known:

\[
\boxed{
v=Z\frac{RT}{P}
}
\]

or

\[
\boxed{v=Zv_{\mathrm{ideal}}}
\]

% ------------------------------------------------------------
\textbf{\\ Important Trends}

At very low reduced pressure:
\[
\boxed{P_R\ll1}
\]

\[
\boxed{Z\rightarrow1}
\]

Hence:
\[
\boxed{
P\rightarrow0
\Rightarrow
\text{real-gas behavior}\rightarrow\text{ideal-gas behavior}
}
\]

At sufficiently high reduced temperature:
\[
\boxed{T_R>2}
\]

ideal-gas behavior is generally a good approximation except at very high reduced pressures.

Near the critical point:
\[
\boxed{|Z-1|\ \text{can become large}}
\]

Thus real-gas deviation is especially important near the critical state.

% ------------------------------------------------------------
\textbf{\\ Pseudo-Reduced Specific Volume}

\[
\boxed{
v_R=
\frac{v_{\mathrm{actual}}}
{RT_{\mathrm{cr}}/P_{\mathrm{cr}}}
}
\]

or

\[
\boxed{
v_R=
\frac{vP_{\mathrm{cr}}}{RT_{\mathrm{cr}}}
}
\]

Used when $P,v$ or $T,v$ are known.

% ------------------------------------------------------------
\textbf{\\ Determining $v$ from $P,T$}

\begin{enumerate}[leftmargin=*]
    \item Obtain $P_{\mathrm{cr}},T_{\mathrm{cr}}$.
    \item Calculate
    \[
    P_R=\frac{P}{P_{\mathrm{cr}}},
    \qquad
    T_R=\frac{T}{T_{\mathrm{cr}}}.
    \]
    \item Obtain $Z$ from the generalized compressibility chart.
    \item Calculate
    \[
    \boxed{v=Z\frac{RT}{P}}.
    \]
\end{enumerate}

\textbf{\\ Determining $P$ from $T,v$}

\[
v_R=\frac{vP_{\mathrm{cr}}}{RT_{\mathrm{cr}}}
\]

\[
T_R=\frac{T}{T_{\mathrm{cr}}}
\]

Use $v_R,T_R$ on the generalized chart to obtain $P_R$, then:

\[
\boxed{
P=P_RP_{\mathrm{cr}}
}
\]

\textbf{\\ Accuracy}

The generalized chart is approximate; accuracy depends on chart resolution, reading precision, and the extent to which the substance obeys the principle of corresponding states.

% ============================================================
\mysection{3.8}{Other Equations of State}{3.8}

\textbf{\\ General Point}

These equations provide improved real-gas representation over a wider region than the ideal-gas equation.

\[
\boxed{
\text{Applicable to gas phase}
}
\]

They should \textbf{not} be used for liquids or liquid--vapor mixtures.

% ------------------------------------------------------------
\textbf{\\ 1. van der Waals Equation}

\[
\boxed{
\left(P+\frac{a}{v^2}\right)(v-b)=RT
}
\]

Meaning:
\begin{itemize}[leftmargin=*]
    \item $\dfrac{a}{v^2}$ accounts for \textbf{intermolecular attraction}.
    \item $b$ accounts for \textbf{molecular volume}.
\end{itemize}

Critical-point conditions:

\[
\boxed{
\left(\frac{\partial P}{\partial v}\right)_{T=T_{\mathrm{cr}}}=0
}
\]

\[
\boxed{
\left(\frac{\partial^2P}{\partial v^2}\right)_{T=T_{\mathrm{cr}}}=0
}
\]

Constants:

\[
\boxed{
a=\frac{27R^2T_{\mathrm{cr}}^2}{64P_{\mathrm{cr}}}
}
\]

\[
\boxed{
b=\frac{RT_{\mathrm{cr}}}{8P_{\mathrm{cr}}}
}
\]

For molar form, use $\bar v$ and $R_u$.

\textbf{Limitation}: conceptually useful and simple, but accuracy is often inadequate compared with more sophisticated equations.

% ------------------------------------------------------------
\textbf{\\ 2. Beattie--Bridgeman Equation}

\[
\boxed{
P=
\frac{R_uT}{\bar v^{\,2}}
\left(1-\frac{c}{\bar v T^3}\right)
(\bar v+B)
-\frac{A}{\bar v^{\,2}}
}
\]

where:

\[
\boxed{
A=A_0\left(1-\frac{a}{\bar v}\right)
}
\]

\[
\boxed{
B=B_0\left(1-\frac{b}{\bar v}\right)
}
\]

Constants are experimentally determined and tabulated.

Applicability:

\[
\boxed{
\rho\leq0.8\rho_{\mathrm{cr}}
}
\]

More accurate than van der Waals over its useful density range.

% ------------------------------------------------------------
\textbf{\\ 3. Benedict--Webb--Rubin Equation}

\[
\boxed{
\begin{aligned}
P={}&
\frac{R_uT}{\bar v}
+
\left(
B_0R_uT-A_0-\frac{C_0}{T^2}
\right)\frac{1}{\bar v^{\,2}}
\\
&+
\frac{bR_uT-a}{\bar v^{\,3}}
+
\frac{a\alpha}{\bar v^{\,6}}
+
\frac{c}{\bar v^{\,3}T^2}
\left(
1+\frac{\gamma}{\bar v^{\,2}}
\right)
e^{-\gamma/\bar v^{\,2}}
\end{aligned}
}
\]

Constants:
\[
\boxed{
A_0,\ B_0,\ C_0,\ a,\ b,\ c,\ \alpha,\ \gamma
}
\]

determined experimentally.

Applicability:

\[
\boxed{
\rho\leq2.5\rho_{\mathrm{cr}}
}
\]

Handles substantially higher gas densities than Beattie--Bridgeman.

% ------------------------------------------------------------
\textbf{\\ Strobridge Extension}

Extension of Benedict--Webb--Rubin with more constants and greater flexibility.

Main use:
\[
\boxed{\text{digital computer calculations}}
\]

% ------------------------------------------------------------
\textbf{\\ 4. Virial Equation}

\[
\boxed{
P=
\frac{RT}{v}
+
\frac{a(T)}{v^2}
+
\frac{b(T)}{v^3}
+
\frac{c(T)}{v^4}
+
\frac{d(T)}{v^5}
+\cdots
}
\]

Coefficients:
\[
\boxed{
a(T),\ b(T),\ c(T),\ d(T),\ldots
}
\]

are \textbf{virial coefficients}, functions of temperature alone.

They may be obtained experimentally or through statistical mechanics.

More terms $\Rightarrow$ potentially greater accuracy.

At low pressure:

\[
\boxed{P\rightarrow0}
\]

\[
\boxed{Pv\rightarrow RT}
\]

\[
\boxed{Z\rightarrow1}
\]

Thus the virial equation is consistent with the ideal-gas limit.

% ------------------------------------------------------------
\textbf{\\ Comparison}

\begin{center}
\begin{tabular}{|c|c|c|c|}
\hline
\textbf{Equation} & \textbf{Constants} & \textbf{Accuracy} & \textbf{Use} \\
\hline
van der Waals & Few & Limited & Simple real-gas modeling \\
\hline
Beattie--Bridgeman & Several & Reasonable & Wider gas-density range \\
\hline
Benedict--Webb--Rubin & More & High & Accurate real-gas calculations \\
\hline
Strobridge & Many & High & Computer-oriented calculations \\
\hline
Virial & Variable & Depends on terms & Series-based representation \\
\hline
\end{tabular}
\end{center}

\textbf{\\ Complexity vs Accuracy}

\[
\boxed{
\text{Greater complexity}
\quad\Longleftrightarrow\quad
\text{generally greater accuracy and wider applicability}
}
\]

\begin{itemize}[leftmargin=*]
    \item \textbf{van der Waals}: simplest; two constants; limited accuracy.
    \item \textbf{Beattie--Bridgeman}: more complex; experimentally determined constants.
    \item \textbf{Benedict--Webb--Rubin}: more constants; higher accuracy; wider density range.
    \item \textbf{Virial}: series form; accuracy increases with retained terms.
\end{itemize}

\textbf{\\ Calculation Consideration}

These equations are generally \textbf{implicit in specific volume}; therefore solving for $v$ may require \textbf{trial-and-error}.

Complex equations are particularly suitable for \textbf{digital computers}, while property tables or simpler equations are often more convenient for hand calculations.

\textbf{\\ Important Limitation}

\[
\boxed{\text{Do not use these gas-phase equations for liquids}}
\]

\[
\boxed{\text{Do not use these gas-phase equations for liquid--vapor mixtures}}
\]

Use property tables or other appropriate relations for those regions.

\end{document}